% Downstream of the approximation convention owned by TNLean #8788.
% Checked locally against the pinned, unmerged model dependency.
% Register after dependency integration; deliberately not in content.tex yet.
\section{The finite-size approximation conclusion}
\label{sec:peps_exact_small_size_approximation}

\begin{theorem}[Approximation on all smaller positive squares]
    \label{thm:peps_exact_small_size_approximation}
    \lean{TNLean.PEPS.Approximation.hasPEPSApproximation_of_small_size}
    \leanok
    \uses{thm:peps_exact_small_sizes}
    Fix integers $q\ge1$ and $L_0\ge0$, and real numbers $C$ and $c\ge0$.
    Set $C'=\max\{C,q^{L_0^2}\}$. For every integer $0<L\le L_0$ and every
    unit vector $\Omega$ on the open $L\times L$ square, there is a nonzero
    PEPS $\Phi$ on that square with $1\le D_e\le C'L^c$ on every edge and a
    real phase $\theta$ such that
    $\|\|\Phi\|^{-1}\Phi-e^{i\theta}\Omega\|\le L^{-1}$.
    This is the finite-size part of~\cite[assembly, lines 203--214]{PolynomialPEPS2026}.
\end{theorem}
\begin{proof}
    \leanok
    \uses{thm:peps_exact_small_sizes}
    Choose the exact representation from Theorem~\ref{thm:peps_exact_small_sizes}.
    Equality of its coefficients with those of $\Omega$ gives $\Phi=\Omega$.
    Thus $\|\Phi\|=1$, so $\Phi\ne0$. Taking $\theta=0$ makes the normalized
    error zero, which is at most $L^{-1}$.
\end{proof}

% The uniform-quantifier wrapper is checked locally at revision 6b194ffc9.
\begin{theorem}[A uniform positive prefactor for a finite range]
    \label{thm:peps_uniform_small_size_approximation}
    \lean{TNLean.PEPS.Approximation.exists_uniform_small_size_approximation}
    \leanok
    \uses{thm:peps_exact_small_size_approximation}
    Fix integers $q\ge1$ and $L_0\ge0$, and real numbers $C$ and $c\ge0$.
    There is a constant $C'>0$ with $C'\ge C$ such that, for every integer
    $2\le L<L_0$ and every unit vector $\Omega$ on the open $L\times L$ square,
    there is a nonzero PEPS $\Phi$ on that square with
    $1\le D_e\le C'L^c$ on every edge and a real phase $\theta$ satisfying
    $\|\|\Phi\|^{-1}\Phi-e^{i\theta}\Omega\|\le L^{-1}$.
    The constant is chosen before $L$ and $\Omega$, as required for the finite-size
    argument in~\cite[assembly, lines 203--214]{PolynomialPEPS2026}.
\end{theorem}
\begin{proof}
    \leanok
    \uses{thm:peps_exact_small_size_approximation}
    Take $C'=\max\{C,q^{L_0^2}\}$. Since $q\ge1$, the second entry of the
    maximum is positive, so $C'>0$ and $C'\ge C$.
    Apply Theorem~\ref{thm:peps_exact_small_size_approximation} at each $L$.
\end{proof}
